% Generated draft: requires mathematical and source review.
\documentclass[11pt]{article}
\usepackage[margin=25mm]{geometry}
\usepackage{fontspec}
\setmainfont{Latin Modern Roman}
\usepackage{amsmath,amssymb,mathtools,unicode-math}
\setmathfont{Latin Modern Math}
\usepackage{enumitem,needspace,xurl,hyperref}
\setlength{\parindent}{0pt}
\setlength{\parskip}{5pt}
\setlength{\emergencystretch}{3em}
\newcommand{\N}{\mathbb{N}}
\newcommand{\Z}{\mathbb{Z}}
\newcommand{\Q}{\mathbb{Q}}
\newcommand{\R}{\mathbb{R}}
\newcommand{\C}{\mathbb{C}}
\newcommand{\F}{\mathbb{F}}
\newcommand{\A}{\mathbb{A}}
\newcommand{\PP}{\mathbb{P}}
\newcommand{\E}{\mathbb{E}}
\newcommand{\eps}{\varepsilon}
\newcommand{\dd}{\,\mathrm{d}}
\begin{document}
\section*{UnsolvedMath ID: 800881}
\textbf{Source problem code:} catalogue entry
\subsection{Definitions and mathematical statement}
For notation, let \(n\geq 1\) and let \(\Omega\subset\C^n\) be a domain. This is a conventional local framework for interpreting the question; the supplied excerpt does not specify the dimension or domain. Write \(\operatorname{PSH}(\Omega)\) for the plurisubharmonic functions on \(\Omega\): upper semicontinuous functions whose restrictions to complex lines are subharmonic or identically \(-\infty\). The condition \(u\in L^\infty_{\mathrm{loc}}(\Omega)\) means that \(u\) is essentially bounded on every relatively compact open subset of \(\Omega\).

Fix the normalization \(dd^c=i\partial\bar\partial\). For locally bounded plurisubharmonic \(u\), interpret \((dd^c u)^n\) as the Bedford--Taylor complex Monge--Amp\`ere measure. Define its range by
\[
\mathcal R(\Omega)=\left\{\mu:\ \mu\text{ is a positive Radon measure on }\Omega,\quad \mu=(dd^c u)^n\text{ for some }u\in\operatorname{PSH}(\Omega)\cap L^\infty_{\mathrm{loc}}(\Omega)\right\}.
\]
The normalization is an explicit convention for this draft, not one specified in the supplied excerpt.

\textbf{Question attributed to Bedford, as recalled in Question 12.} Can membership in \(\mathcal R(\Omega)\) be characterized by a uniform integrability condition?

The supplied formulation does not specify the family of functions or measures to which uniform integrability should apply, a reference measure or capacity, or whether the condition is local or global. Thus it does not supply a definite necessary-and-sufficient condition that can be written as an equivalence. It also gives no boundary conditions or total-mass restriction. These omissions remain unresolved rather than being added as hypotheses.

The surrounding discussion concerns bounded and continuous plurisubharmonic functions, and also mentions quasi-plurisubharmonic functions. The displayed Bedford question specifically concerns \(\operatorname{PSH}\cap L^\infty_{\mathrm{loc}}\); continuity, global boundedness, and the quasi-plurisubharmonic setting are not imposed in that question.
\subsection{Short English statement}
Bedford asks whether the measures obtained as complex Monge--Amp\`ere measures of locally bounded plurisubharmonic functions admit a characterization by uniform integrability. The excerpt leaves the proposed integrability condition and several ambient specifications unstated, so this remains a qualified question rather than a fully specified equivalence.
\subsection{Sources}
\begin{itemize}
\item Local UnsolvedMath registry: permanent ID \(800881\), source problem code \texttt{catalogue entry}, title \emph{Open Problems in Pluripotential Theory --- Question 12}. The supplied registry labels the problem \emph{open}; this is a supplied status, not an independently verified mathematical or literature assessment.
\item Dinew, Guedj, and Zeriahi, \emph{Open Problems in Pluripotential Theory} (2015), Question 12. Supplied source URL: \url{https://arxiv.org/abs/1511.00705}. The provenance records access on 2026-10-16 and extraction by \texttt{arxiv\_latex}; these are supplied provenance statements, not a claim of access or verification for this draft. The registry lists Dinew, Guedj, and Zeriahi as proposers in 2015, while the excerpt attributes this particular question to Bedford.
\item Reference label [Bed93], p.~70: the excerpt identifies this as the earlier source of Bedford's question. Its full bibliographic details are not supplied.
\item Reference labels [BT76] and [BT82]: the excerpt attributes the original breakthrough concerning bounded and continuous plurisubharmonic functions to Bedford and Taylor. Full bibliographic details are not supplied.
\item Reference label [Kol98]: the excerpt cites this for rather sharp sufficient conditions while reporting that the range was not characterized. This is the supplied source's assessment; full bibliographic details and independent verification are absent.
\end{itemize}
\end{document}
